%% ============================================================
%% PARTE VI — GUIA DE OBSERVAÇÃO E ENCAMINHAMENTO (R201.4)
%% Rastreio escolar por sinais de alerta + condições com CID
%% Uso exclusivo do professor e do psicólogo institucional
%% ============================================================

\chapter{Guia de Observação e Encaminhamento}

\section{Aviso Legal e Ético --- Leia Antes de Usar}

\begin{avaliacao}
\textbf{Este guia NÃO diagnostica.} As fichas deste capítulo são instrumentos
de \textbf{observação escolar}: o professor registra sinais de alerta e
\textbf{encaminha} a criança aos profissionais competentes. O diagnóstico de
dislexia, TEA, TDAH ou qualquer outra condição clínica é ato privativo de
psicólogo (Lei nº 4.119/1962; Resolução CFP nº 009/2018), psiquiatra da
infância, neurologista pediátrico ou fonoaudiólogo, conforme o caso. Nenhuma
ficha deste livro gera rótulo, laudo ou classificação de aluno.
\end{avaliacao}

\section{Quem Faz o Quê no Rastreio}

\begin{center}
\resizebox{\textwidth}{!}{%
\begin{tabular}{|l|l|l|}
\hline
\textbf{Profissional} & \textbf{Papel no rastreio} & \textbf{O que NÃO pode fazer} \\
\hline
Professor & Observar, registrar, adaptar, encaminhar & Diagnosticar, rotular, aplicar teste psicológico \\
\hline
Professor & Observar, registrar, adaptar, encaminhar & Diagnosticar, rotular, aplicar teste psicológico \\
\hline
Psicólogo institucional & Triar, orientar equipe, apoiar o professor & Liberar laudo diagnóstico sem avaliação clínica completa \\
\hline
Psicólogo/Neuropsicólogo & Avaliação cognitiva com instrumentos validados & Deixar de exigir anamnese e consentimento \\
\hline
Fonoaudiólogo & Avaliar linguagem, fala, consciência fonológica & Diagnosticar TEA/TDAH isoladamente \\
\hline
Médico (neuro/pediatra/psiquiatra) & Investigação clínica e etiológica & Diagnosticar só por relato escolar, sem exames \\
\hline
\end{tabular}}
\end{center}

\begin{avaliacao}
\textbf{Integração com a versão profissional da coleção:} os instrumentos de
\textbf{sondagem inicial minuciosa} (11 domínios) e as \textbf{30 fichas de
rastreio sistemático e investigativo} desta coleção estão no
\textbf{Volume Profissional}, de uso restrito a psicólogo, psicopedagogo e
neuropsicopedagogo (avaliação do aluno e orientação de professores). Este
capítulo contém apenas orientações de observação escolar e encaminhamento;
o material do aluno permanece sem avaliação clínica.
\end{avaliacao}

\section{Base Legal}

\begin{itemize}
  \item \textbf{Lei nº 4.119/1962} e \textbf{Resolução CFP nº 009/2018}: uso
        de testes psicológicos é privativo de psicólogo.
  \item \textbf{ECA (Lei nº 8.069/1990)}: proteção integral da criança;
        vedação de tratamento vexatório ou estigmatizante.
  \item \textbf{LGPD (Lei nº 13.709/2018), Art. 11}: dado de saúde é dado
        sensível — exige consentimento específico dos responsáveis, finalidade
        determinada e mínimo necessário. Fichas de observação com sinais
        clínicos são dados sensíveis e devem ficar em local restrito.
  \item \textbf{Lei nº 13.146/2015 (Estatuto da Pessoa com Deficiência)}:
        direito à educação inclusiva com adaptações.
  \item \textbf{Resolução CNE/CEB nº 7/2010, Arts. 29 e 37}: atendimento à
        diversidade de ritmos e condições de aprendizagem.
\end{itemize}

%% ============================================================
\section{Fichas de Observação Estruturada por Domínio}

\subsection{Como Usar as Fichas}

\begin{enumerate}[leftmargin=2cm]
  \item Observe por pelo menos \textbf{4 semanas} em contextos variados
        (sala, pátio, atividade dirigida, atividade livre).
  \item Marque \textbf{S} (sempre), \textbf{F} (frequentemente),
        \textbf{R} (raramente) ou \textbf{N} (nunca).
  \item \textbf{Sinal de alerta}: comportamento marcado S ou F que persiste
        após as acomodações do livro (rotas 1A--1D, pausas, pistas) e que
        destoa claramente da turma.
  \item Dois ou mais sinais de alerta no mesmo domínio $\rightarrow$
        \textbf{encaminhar}: registre por escrito, reúna-se com a família e
        acione o psicólogo institucional. Nunca comunique hipótese diagnóstica
        à família: comunique \textbf{fatos observados} e \textbf{necessidades}.
\end{enumerate}

\subsection{D1 --- Linguagem Oral e Comunicação}

\begin{center}
\resizebox{\textwidth}{!}{%
\begin{tabular}{|l|c|c|c|c|}
\hline
\textbf{Indicador observável} & S & F & R & N \\
\hline
Pronuncia a maioria das palavras de forma compreensível para estranhos & & & & \\
\hline
Tem vocabulário compatível com a idade para relatar fatos & & & & \\
\hline
Compreende instruções simples sem repetição & & & & \\
\hline
Organiza frases com verbo e sujeito & & & & \\
\hline
Mantém o tema da conversa & & & & \\
\hline
\end{tabular}}
\end{center}

\caixapista{Alerta se: trocas de fala persistentes após 5 anos (F80.0), frases
curtas ou ausência de narrativa (F80.1/F80.2), gagueira com tensão (F98.5).
Profissional de referência: fonoaudiólogo.}

\subsection{D2 --- Consciência Fonológica e Leitura}

\begin{center}
\resizebox{\textwidth}{!}{%
\begin{tabular}{|l|c|c|c|c|}
\hline
\textbf{Indicador observável} & S & F & R & N \\
\hline
Percebe rimas e palavras com o mesmo início & & & & \\
\hline
Segmenta palavras em sílabas batendo palmas & & & & \\
\hline
Associa letra--som com precisão nas folhas do livro & & & & \\
\hline
Funde sílabas para ler palavras novas (K3) & & & & \\
\hline
Memoriza palavras lidas com frequência & & & & \\
\hline
Troca letras espelhadas (b/d, p/q) além do esperado para a etapa & & & & \\
\hline
\end{tabular}}
\end{center}

\caixapista{Alerta se: após instrução fônica sistemática (este livro) e Rota
1A/1B, persiste troca persistente, leitura adivinhada ou extrema lentidão.
Referência: psicopedagogo + fonoaudiólogo; instrumentos do especialista:
CONFIAS, TDE-II, PCFO.}

\subsection{D3 --- Escrita e Motricidade Fina}

\begin{center}
\resizebox{\textwidth}{!}{%
\begin{tabular}{|l|c|c|c|c|}
\hline
\textbf{Indicador observável} & S & F & R & N \\
\hline
Segura o lápis com preensão adequada & & & & \\
\hline
Traça as 4 formas da letra respeitando direção & & & & \\
\hline
Escreve dentro da linha & & & & \\
\hline
Omissões/trocas de letras diminuem com as missões de consolidação & & & & \\
\hline
Realiza as missões motoras sem frustração intensa & & & & \\
\hline
\end{tabular}}
\end{center}

\caixapista{Alerta se: disgrafia persistente (F81.8), coordenação motora
globalmente atrasada (F82). Referência: terapeuta ocupacional/psicomotricista
+ neurologista pediátrico.}

\subsection{D4 --- Atenção, Funções Executivas e Autorregulação}

\begin{center}
\resizebox{\textwidth}{!}{%
\begin{tabular}{|l|c|c|c|c|}
\hline
\textbf{Indicador observável} & S & F & R & N \\
\hline
Inicia a tarefa após uma instrução & & & & \\
\hline
Sustenta a atenção por 10--15 minutos em atividade do livro & & & & \\
\hline
Conclui a plancheta sem abandonar & & & & \\
\hline
Espera a vez e inibe respostas impulsivas & & & & \\
\hline
Usa a caixa de Pausa quando precisa (autorregulação) & & & & \\
\hline
Organiza o material & & & & \\
\hline
\end{tabular}}
\end{center}

\caixapista{Alerta se: desatenção/impulsividade/hiperatividade intensas,
pervasivas (casa e escola) e por mais de 6 meses $\rightarrow$ avaliação para
TDAH (F90.0/F90.1/F90.2). Referência: neuropsicólogo + psiquiatra infantil.
Instrumentos do especialista: SNAP-IV, Escala Wechsler (WISC-V).}

\subsection{D5 --- Comportamento Social e Emocional}

\begin{center}
\resizebox{\textwidth}{!}{%
\begin{tabular}{|l|c|c|c|c|}
\hline
\textbf{Indicador observável} & S & F & R & N \\
\hline
Faz contato visual e responde ao nome & & & & \\
\hline
Inicia e responde a interações com colegas & & & & \\
\hline
Tolera mudanças de rotina com apoio visual & & & & \\
\hline
Expressa emoções de forma proporcional à situação & & & & \\
\hline
Aceita contato físico e aproximação & & & & \\
\hline
\end{tabular}}
\end{center}

\caixapista{Alerta se: pouca interação social recíproca, interesses restritos,
estereotipias (F84.0/F84.1); mutismo em contexto específico (F94.0); ansiedade
intensa de separação (F93.0). Referência: neuropsicólogo/psiquiatra infantil.
Instrumentos do especialista: o M-CHAT-R/F aplica-se somente a 16--30 meses
(não se usa nesta faixa escolar); para 6--10 anos, triagem com SRS-2 e
confirmação com ADOS-2/ADI-R, CARS-2 — aplicados por profissional habilitado.}

\subsection{D6 --- Memória e Ritmo de Aprendizagem}

\begin{center}
\resizebox{\textwidth}{!}{%
\begin{tabular}{|l|c|c|c|c|}
\hline
\textbf{Indicador observável} & S & F & R & N \\
\hline
Retém o conteúdo da missão anterior na revisão cumulativa & & & & \\
\hline
Generaliza famílias silábicas para palavras novas & & & & \\
\hline
Necessita de repetição dentro do esperado para a rota & & & & \\
\hline
Esquece instruções de múltiplas etapas & & & & \\
\hline
\end{tabular}}
\end{center}

\caixapista{Alerta se: defasagem global persistente em todos os domínios com
apoio adequado $\rightarrow$ investigar deficiência intelectual (F70--F79) ou
causas clínicas. Referência: psicólogo (avaliação cognitiva) + médico.}

%% ============================================================
\section{Tabela de Condições e CIDs (Referência Institucional)}

A tabela reúne as condições mais relevantes no contexto do 1º ano. Os códigos
são do CID-10, vigente na documentação brasileira; as equivalências CID-11
aparecem entre colchetes quando relevantes. A tabela serve para \textbf{registro
institucional e comunicação com a rede de saúde} --- nunca para classificar o
aluno.

\begin{center}
\resizebox{\textwidth}{!}{%
\begin{tabular}{|l|p{5.5cm}|p{4.5cm}|p{4cm}|}
\hline
\textbf{CID-10} & \textbf{Condição} & \textbf{Profissional de avaliação} & \textbf{Acomodação escolar principal} \\
\hline
F81.0 & Dislexia (transtorno específico de leitura) & Neuropsicólogo + fonoaudiólogo & Método fônico sistemático, tempo extra, sem exposição de leitura oral \\
\hline
F81.1 & Disortografia (soletração) & Neuropsicólogo + psicopedagogo & Não punir erros ortográficos; valorizar conteúdo \\
\hline
F81.2 & Discalculia & Neuropsicólogo + psicopedagogo & Material concreto, foco no raciocínio \\
\hline
F81.3 & Transtorno misto de habilidades escolares & Neuropsicólogo & Acomodações combinadas \\
\hline
F81.8 & Disgrafia e outros transtornos escolares & Neuropsicólogo + T.O. & Digitação/alternativas à escrita cursiva \\
\hline
F82 & Transtorno do desenvolvimento da coordenação & T.O./psicomotricista + neuro & Missões motoras ampliadas, sem pressa \\
\hline
F80.0 & Transtorno fonológico (articulação) & Fonoaudiólogo & Apoio articulatório (pistas da boca) \\
\hline
F80.1 & Transtorno da linguagem expressiva & Fonoaudiólogo & Resposta por gesto/desenho aceita \\
\hline
F80.2 & Transtorno da linguagem receptiva & Fonoaudiólogo & Instruções curtas + apoio visual \\
\hline
F98.5 & Gagueira & Fonoaudiólogo & Nunca completar a frase; tempo total de fala \\
\hline
F84.0 & Autismo infantil (TEA) [6A02] & Neuropsicólogo/psiquiatra infantil & Rotina visual, antecipação, interesses como ponte \\
\hline
F84.1 & Autismo atípico & Idem & Idem \\
\hline
F84.5 & Síndrome de Asperger (CID-10; integrada ao TEA na CID-11) & Idem & Rotina + apoio social explícito \\
\hline
F84.2 & Síndrome de Rett & Neurologista pediátrico & Plano individualizado multidisciplinar \\
\hline
F90.0 & TDAH --- apresentação desatenta [6A05] & Neuropsicólogo + psiquiatra & Uma instrução por vez, estímulos reduzidos \\
\hline
F90.1 & TDAH --- hiperativo-impulsivo & Idem & Pausas de movimento legítimas, nunca punitivas \\
\hline
F90.2 & TDAH --- apresentação combinada & Idem & Combinar as duas estratégias \\
\hline
F91.3 & Transtorno desafiador opositivo (TOD) & Psiquiatra infantil + psicólogo & Regras claras, reforço positivo, evitar confronto \\
\hline
F91 & Transtorno de conduta & Psiquiatra infantil & Limites consistentes + acompanhamento \\
\hline
F95.2 & Síndrome de Tourette & Neurologista & Tiques ignorados; pausas sem estigma \\
\hline
F93.0 & Ansiedade de separação & Psicólogo + psiquiatra & Rotina de acolhimento previsível \\
\hline
F94.0 & Mutismo seletivo & Psicólogo + fonoaudiólogo & Nunca forçar a fala; aceitar resposta escrita \\
\hline
F94.1 & Transtorno reativo de vinculação & Psicólogo + equipe & Vínculo estável com adulto de referência \\
\hline
F42 & Transtorno obsessivo-compulsivo (início precoce) & Psiquiatra infantil & Flexibilizar rituais sem confrontar \\
\hline
F41 & Ansiedade generalizada & Psicólogo + psiquiatra & Reduzir antecipação negativa; pausas \\
\hline
F32/F33 & Depressão infantil & Psiquiatra infantil + psicólogo & Atividades curtas, acolhimento, ativar vínculo \\
\hline
F31 & Transtorno bipolar (início precoce) & Psiquiatra infantil & Comunicação estreita com a equipe de saúde \\
\hline
F70--F79 & Deficiência intelectual & Psicólogo (cognitivo) + médico & Currículo adaptado ao nível real, passos menores \\
\hline
H90--H91 & Deficiência auditiva & Otorrino + fonoaudiólogo & Posicionamento frontal, apoio visual, libras se necessário \\
\hline
H54 & Baixa visão/cegueira & Oftalmologista & Fonte ampliada (este livro já usa), contraste, braille \\
\hline
G40 & Epilepsia & Neurologista & Plano de crise conhecido pela equipe \\
\hline
G80 & Paralisia cerebral & Neurologista + equipe & Adaptação de acesso à escrita \\
\hline
Q90 & Síndrome de Down & Equipe multidisciplinar & Passos menores, reforço positivo, autonomia \\
\hline
Q99.2 & Síndrome do X frágil & Genética + equipe & Rotina estável, reduzir hiperestimulação \\
\hline
Q87.1 & Síndrome de Prader-Willi & Equipe multidisciplinar & Supervisão alimentar, rotina clara \\
\hline
F98.0 & Enurese & Pediatra & Discrição, sem punição \\
\hline
F98.1 & Encoprese & Pediatra + psicólogo & Discrição, sem humilhação \\
\hline
H93.2 & Transtorno do processamento auditivo central & Otorrino + fonoaudiólogo & Sentar à frente, falar de frente, reduzir ruído \\
\hline
\end{tabular}}
\end{center}

\caixapausa{Esta tabela é longa de propósito: é material de consulta do
professor e do psicólogo institucional, não página de atividade.}

%% ============================================================
\section{Fluxo de Encaminhamento}

\begin{center}
\begin{tikzpicture}[
  etapa/.style={rectangle, draw=primary, fill=primary!10, align=center,
                minimum width=5.5cm, minimum height=1.1cm, font=\small},
  seta/.style={->, thick, color=accent}
]
  \node[etapa] (e1) {1. Observar 4 semanas\\com as fichas D1--D6};
  \node[etapa, below=0.6cm of e1] (e2) {2. Aplicar acomodações\\rotas 1A--1D + pausas + pistas};
  \node[etapa, below=0.6cm of e2] (e3) {3. Persistiu? Registrar fatos\\+ reunir com a família};
  \node[etapa, below=0.6cm of e3] (e4) {4. Acionar psicólogo\\institucional (triagem)};
  \node[etapa, below=0.6cm of e4] (e5) {5. Encaminhar à rede:\\fono, neuropsicólogo, médico};
  \node[etapa, below=0.6cm of e5] (e6) {6. Receber laudo?\\Adaptar o PEI + acolher};
  \draw[seta] (e1) -- (e2);
  \draw[seta] (e2) -- (e3);
  \draw[seta] (e3) -- (e4);
  \draw[seta] (e4) -- (e5);
  \draw[seta] (e5) -- (e6);
\end{tikzpicture}
\end{center}

\section{Modelo de Registro de Encaminhamento}

\begin{center}
\resizebox{\textwidth}{!}{%
\begin{tabular}{|l|p{9cm}|}
\hline
\textbf{Campo} & \textbf{Registro} \\
\hline
Aluno (iniciais apenas) & \underline{\hspace{6cm}} \\
\hline
Domínios com sinais de alerta (D1--D6) & \underline{\hspace{6cm}} \\
\hline
Acomodações já aplicadas & \underline{\hspace{6cm}} \\
\hline
Fatos observados (sem hipóteses diagnósticas) & \underline{\hspace{6cm}} \\
\hline
Data do contato com a família & \underline{\hspace{4cm}} \\
\hline
Encaminhamento sugerido & \underline{\hspace{6cm}} \\
\hline
\end{tabular}}
\end{center}

\caixapista{Regra de ouro do registro: descreva o que a criança FAZ, nunca o
que ela ``é''. ``Não lê palavras com sílabas diretas após 8 semanas de Rota
1B'' é fato. ``É disléxico'' é rótulo --- e rótulo não se escreve em ficha
escolar.}

\checklistdominio{Protocolo de observação aplicado com ética: fatos, não rótulos; encaminhamento, não diagnóstico}
